% id: 0625
% book: RPM 중학 3-1 · 05 다항식의 인수분해
% section: 시험에 꼭 나오는 문제
% figure: crop:fig-0625.png
% answer: ③
% answer_source: 답지
% variant_level: 0 (원본 전사)
% 이 파일은 build.mjs 가 items.json 에서 만든다. 고칠 때는 items.json 을 고친다.
\smash{\raisebox{1.5mm}{\dmkichul{RPM 중학 3-1 0625번 · 88쪽}}}
다음 그림과 같이 한 변의 길이가 $k$인 정사각형에서 한 변의 길이가 $3$인 정사각형을 잘라 내고 남은 도형을 반으로 잘라 붙여서 새로운 직사각형을 만들었다. 두 도형의 넓이가 같음을 이용하여 설명할 수 있는 인수분해 공식은?
\par\smallskip\begin{center}\includegraphics[width=149.4pt]{figures/fig-0625.png}\end{center}
\begin{choicesv}{$x^2+2xy+y^2=(x+y)^2$}{$x^2-2xy+y^2=(x-y)^2$}{$x^2-y^2=(x+y)(x-y)$}{$x^2+(a+b)x+ab=(x+a)(x+b)$}{$acx^2+(ad+bc)x+bd=(ax+b)(cx+d)$}\end{choicesv}
